\documentclass[11pt,a4paper]{article}
\usepackage[margin=23mm]{geometry}
\usepackage{amsmath,amssymb,amsthm,booktabs,hyperref}
\newtheorem{theorem}{Theorem}[section]
\newtheorem{lemma}[theorem]{Lemma}
\newtheorem{proposition}[theorem]{Proposition}
\newtheorem{corollary}[theorem]{Corollary}
\newcommand{\F}{\mathbb F}
\newcommand{\GL}{\operatorname{GL}}
\newcommand{\PSU}{\operatorname{PSU}}
\newcommand{\PGU}{\operatorname{PGU}}
\newcommand{\PGammaU}{\operatorname{P\Gamma U}}
\newcommand{\End}{\operatorname{End}}
\newcommand{\Aut}{\operatorname{Aut}}
\title{Local actions, invariant pencils and scalar quotients}
\author{Human Compression project: private unreviewed draft}
\date{}
\begin{document}\maketitle
This common paper develops the local-action arguments used in a complete
corrected recovery of Buekenhout--Dehaye--Leemans~\cite{BDL}. The complete application is written separately and remains subject to distinct
review. Large classification
inputs will remain explicitly attributed; none is inferred from the small
calculations accompanying this draft.

\section{Local actions and preserved partitions}
A finite linear space has a unique line through two distinct points, at
least two points on each line, and at least two lines on each point. The
specified faithful group $G$ is part of the data. Its pencil at $x$ is the
set $\mathcal L(x)$ of lines through $x$. We use local primitivity for
primitivity of both $G_x$ on $\mathcal L(x)$ and $G_L$ on the points of $L$.
For a flag-transitive linear space this is equivalent to the source's
residual weak primitivity: the remaining global weak-primitivity condition
follows from the retained Higman--McLaughlin point-primitivity theorem.
For translation structures below it also follows directly from the proof.

\begin{lemma}[Supergroups and coarsening]\label{coarsening}
A permutation group containing a primitive subgroup is primitive, and one
containing a two-transitive subgroup is two-transitive. Suppose a group
$H$ acts transitively on a finite set $X$ and preserves a partition of $X$
with at least two parts, each of size at least two. Then its action is not
primitive. In particular, this conclusion applies to a pencil partition
induced by a genuinely $H$-invariant coarser incidence structure.
\end{lemma}
\begin{proof}
A block system for the larger group is one for the subgroup; transitivity
of the subgroup rules out extra orbits. Transitivity on ordered distinct
pairs also passes to a supergroup. The partition in the second assertion
is, by hypothesis, a nontrivial invariant equivalence relation, which is
exactly an obstruction to primitivity. Notice that ordering two line
sizes alone does not construct such a relation.
\end{proof}

\begin{proposition}[Translation local actions]\label{translation}
Let $V$ be a finite vector space over $\F_p$, let $\mathcal S$ be a spread
of nonzero proper additive subspaces, and let $H\leq\GL(V)$ preserve
$\mathcal S$ and act transitively on it. The cosets of members of
$\mathcal S$ are the lines of a translation linear space. Put $G=V\rtimes H$.
Then $G$ is flag-transitive. It is locally primitive exactly when $H$ is
primitive on $\mathcal S$ and $H_U$ acts irreducibly on $U$ for a member
$U\in\mathcal S$. It is locally two-transitive exactly when $H$ is
two-transitive on $\mathcal S$ and $H_U$ is transitive on $U\setminus\{0\}$.
In the primitive case $G$ is point-primitive as well.
\end{proposition}
\begin{proof}
Translations are transitive on points and move each affine pencil to the
pencil at zero. The latter pencil is $\mathcal S$, with induced group $H$.
The setwise stabilizer of the line $U$ induces $U\rtimes H_U|_U$ on it.
For a regular translation group, a block through zero is an additive
subgroup: if $b$ lies in the block, translation by $b$ makes the block
meet itself and hence fixes it setwise. Conversely every invariant additive
subgroup gives its coset block system. Thus the affine line action is
primitive exactly when its linear stabilizer is irreducible. It is
two-transitive exactly when the zero stabilizer is transitive on nonzero
vectors. This proves both equivalences, including possible kernels of the
restriction $H_U\to\GL(U)$.

Finally let $W$ be a nonzero $H$-invariant additive subspace of $V$.
Choose $0\ne w\in W$ and the unique spread member $U$ containing it.
The intersection $W\cap U$ is $H_U$-invariant, so irreducibility gives
$U\subseteq W$. Transitivity on the spread gives every member inside
$W$, and these members cover $V\setminus\{0\}$. Hence $W=V$, proving
point-primitivity without a separate global classification.
\end{proof}

\begin{lemma}[A field detected by a commutant]\label{field}
Suppose $N\leq\GL_{\F_p}(V)$ and its commuting algebra is a field
$Q=\End_{\F_pN}(V)$ acting as scalars on $V$. Every linear transformation
normalizing $N$ is semilinear on the $Q$-space $V$.
If $\F_q\subsetneq Q$, $\dim_Q V\geq2$, and $H$ normalizes $N$, then
$H$ is imprimitive on the $\F_q$-lines whenever it acts on those lines:
the $Q$-lines give a nontrivial invariant partition.
\end{lemma}
\begin{proof}
For $h\in N_{\GL_{\F_p}(V)}(N)$, conjugation sends the commuting algebra
of $N$ onto itself. Its restriction to $Q$ is a field automorphism
$\sigma_h$, and $h(av)=\sigma_h(a)h(v)$. Thus $h$ permutes $Q$-lines.
Each $Q$-line contains exactly $(|Q|-1)/(q-1)>1$ distinct $\F_q$-lines,
and there is more than one $Q$-line. Different $Q$-lines are disjoint away
from zero. Lemma~\ref{coarsening} applies.
\end{proof}

For example, for the natural module of $\operatorname{SL}_m(Q)$, $m\geq2$,
the commuting algebra over $\F_p$ is exactly $Q$. Indeed commuting with
all $I+tE_{ij}$ first forces the same additive map on every coordinate
and then forces that map to commute with multiplication by every $t\in Q$.
It is multiplication by a scalar. This supplies the actual invariant
partition in the source's extension-field special-linear branch. The
corresponding symplectic and exceptional module identifications must be
stated at their precise retained scope; the assertion is not justified by
an abstract subgroup name alone.

\begin{proposition}[Circles]\label{circles}
Let $\Omega$ have at least three points and let its lines be all two-element
subsets. A faithful group $G$ is flag-transitive on this linear space if
and only if it is two-transitive on $\Omega$. In that case it is residually
weakly primitive if and only if $G_x$ is primitive on $\Omega\setminus\{x\}$,
and locally two-transitive if and only if $G$ is three-transitive.
If $G$ also preserves a thick linear space on the same point set, its
circle action is not residually weakly primitive.
\end{proposition}
\begin{proof}
Flags $(x,\{x,y\})$ are ordered pairs of distinct points. The pencil at
$x$ identifies with $\Omega\setminus\{x\}$. A line stabilizer interchanges
its two points by two-transitivity, so its two-point action is already
primitive and two-transitive. Global point-primitivity follows from
two-transitivity. This proves both equivalences.

In the final assertion the sets $L\setminus\{x\}$, for thick-space lines
$L$ through $x$, partition $\Omega\setminus\{x\}$. Every part has at least
two elements and there are at least two parts. They form an invariant
partition for $G_x$, so Lemma~\ref{coarsening} applies. In particular the
word ``thick'' cannot be omitted when using this test on $\operatorname{AG}(m,q)$:
for $q=2$ that space is itself the circle.
\end{proof}

\section{Classical pencils and the small affine exclusions}
\begin{proposition}[Classical local actions]\label{classical}
The natural groups containing $\operatorname{PSL}_{m+1}(q)$ on
$\operatorname{PG}(m,q)$, or containing translations and
$\operatorname{SL}_m(q)$ on $\operatorname{AG}(m,q)$, are flag-transitive
and locally two-transitive for $m\ge2$. In the affine plane the induced
special-linear pencil group is $\operatorname{PSL}_2(q)$; it need not
contain $\operatorname{PGL}_2(q)$.
\end{proposition}
\begin{proof}
Special-linear matrices are transitive on nonzero vectors and on ordered
pairs of distinct one-dimensional subspaces: choose independent
representatives, extend them to bases, and rescale a representative or a
complementary vector to correct the determinant. In the projective pencil
at a point, any quotient matrix $A$ extends to the determinant-one matrix
$\operatorname{diag}((\det A)^{-1},A)$. On a projective line, any
$\operatorname{GL}_2(q)$ action extends by its inverse determinant on a
complementary coordinate. Both local projective actions are therefore
two-transitive. The affine pencil is precisely the projective action
on vector lines. Its vector-line stabilizer contains every nonzero
multiplier on that line, compensating the determinant on a complement;
line translations give two-transitivity there. Supergroups preserve these
properties.

For $m=2$ the scalar kernel in $\operatorname{SL}_2(q)$ has order
$\gcd(2,q-1)$, giving exactly $\operatorname{PSL}_2(q)$ on the pencil.
For odd $q$ this has index two in $\operatorname{PGL}_2(q)$.
Thus $\operatorname{ASL}_2(3)$ on $\operatorname{AG}(2,3)$ has pencil
group $A_4$, not $S_4$, despite being locally two-transitive. The
$\operatorname{PGL}_2(q)$ inclusions in source Propositions 3.7 and 3.8
must be replaced by $\operatorname{PSL}_2(q)$. The source's final
special-linear branches do not require this erroneous inclusion.
\end{proof}

\begin{proposition}[Natural symplectic pencils]\label{symplectic}
For even $m\ge4$, the natural $\operatorname{Sp}_m(q)$ is primitive but
not two-transitive on vector lines. Its normalizer over $\F_p$ preserves
perpendicularity, up to field automorphism and nonzero form multiplier.
Consequently every group between it and its normalizer has a locally
primitive, but not locally two-transitive, affine-space action.
\end{proposition}
\begin{proof}
Symplectic basis extension shows that the stabilizer of a vector line
$P$ has three orbits on vector lines: $P$, the remaining points of
$P^\perp$, and those outside $P^\perp$. Write
\[
 v=(q^m-1)/(q-1),\qquad a=(q^{m-1}-1)/(q-1).
\]
A block through $P$ is a union of stabilizer orbits, so its possible
proper nontrivial sizes are $a$ and $q^{m-1}+1$. But
$\gcd(a,v)=1$, and $v/2<q^{m-1}+1<v$; neither size divides $v$.
This proves primitivity, and the two nonempty other orbits disprove
two-transitivity.

The commuting algebra over $\F_p$ is $\F_q$. Indeed commuting with
every transvection $x\mapsto x+t\langle x,w\rangle w$ forces an
endomorphism to preserve each $\F_qw$, the image after subtracting the
identity. An additive map preserving all vector lines is a common scalar:
apply it to $u,v,u+v,u+cv$ for independent $u,v$ and $c\in\F_q$.
Lemma~\ref{field} now makes a normalizer element semilinear over $\F_q$.
After removing its field automorphism, pulling back the alternating form
gives another invariant form. Relative to the original nondegenerate
form its defining operator commutes with the group, and hence is a
scalar. Perpendicularity is preserved even by the full normalizer.
Finally the symplectic group is transitive on nonzero vectors, so a
vector-line stabilizer is transitive on the nonzero vectors of that
line. Proposition~\ref{translation} gives the remaining local assertions.
\end{proof}

\begin{lemma}[A small normal subgroup test]\label{normaltest}
A nontrivial normal subgroup in a faithful primitive permutation group
is transitive. In particular its order is at least the degree.
A faithful primitive group with a nontrivial center is regular of prime
order. These tests apply to the induced group after removal of the
restriction kernel.
\end{lemma}
\begin{proof}
Normal-subgroup orbits are a block system. All singleton orbits would
make the subgroup trivial by faithfulness. A central transitive subgroup
is regular, since an element fixing one point fixes all. Its centralizer
in the symmetric group is its own regular copy, forcing the whole group
to be abelian regular. Primitivity of a regular group is equivalent to
having no proper nontrivial subgroup, hence to prime order.
\end{proof}

Retain Huybrechts' actual representation cases as quoted in source
Proposition 3.4, rather than just abstract isomorphic subgroup names.
The exceptional degree-364 pencil group is $\operatorname{PSL}_2(13)$,
of order $1092$. A point stabilizer has order three and lies properly
in a cyclic split torus of order six: an order-three element diagonalizes
over $\F_{13}$, and the determinant-one diagonal torus modulo scalars
has order six. Thus this action is imprimitive. In the degree-40 cases
$2^4.A$, the normal subgroup of order sixteen has a nontrivial image
of order at most sixteen on directions: the scalar kernel in the
linear $\F_3$-action has order at most two. Such an image cannot be
transitive of degree forty, so Lemma~\ref{normaltest} excludes
primitivity. In the remaining $4.A_5$ or $N.S_5$ cases ($|N|\leq4$),
the specified normal subgroup of order at most four likewise has either
a nontrivial, nontransitive image, or lies entirely in the restriction
kernel. In the latter case the induced group is a quotient of $A_5$
or $S_5$. The former order is not divisible by forty. A transitive
$S_5$ quotient of degree forty must be $S_5$ itself (its other proper
quotients have order at most two); its order-three point stabilizer
lies properly in an $S_3$. Thus this action is imprimitive as well.
This supplies the missing explicit local exclusions, conditional on
the retained representation classification.

\begin{proposition}[The exceptional affine plane]\label{nineplane}
Retain Foulser's ten flag-transitive actions on $\operatorname{AG}(2,9)$
with induced pencil group $A_5$ or $S_5$, each containing the specified
minimal $F=3^4\rtimes2.A_5$ in its natural $\operatorname{SL}_2(9)$
representation. All ten are locally primitive and none is locally
two-transitive.
\end{proposition}
\begin{proof}
In the degree-ten action $A_5$ has point stabilizer $S_3$, which is
maximal. Indeed a proper overgroup could only have order twelve or
thirty; an order-twelve subgroup is $A_4$ and contains no $S_3$, and
simplicity forbids index two. Neither $60$ nor $120$ is divisible by
$10\cdot9$, excluding two-transitivity for either pencil image.

Choose an involution of $A_5$ and a lift $f$ in $2.A_5$. The only
involution in $\operatorname{SL}_2(9)$ is $-I$: a determinant-one
involution in odd characteristic is diagonalizable with both eigenvalues
$-1$, unless it is the identity. Thus $f$ has order four and $f^2=-I$.
The polynomial $X^2+1$ has two distinct roots in $\F_9$, so $f$ fixes
two vector lines. On either line its scalar of order four generates
$\F_9$ over $\F_3$. It acts irreducibly on the additive two-dimensional
line. With line translations this gives a primitive line action;
flag transitivity transfers it to every line. Supergroups preserve both
local primitivity conditions.

Only the cyclic subgroup generated by $f$ is required to fix a direction.
The whole Sylow subgroup need not do so. This avoids the stronger
simultaneous Sylow-fixed-point assertion in the printed proof of
Proposition 3.9 while supplying its claimed conclusion.
\end{proof}

For the other Foulser plane cases, a faithful primitive action of $A_4$
or $S_4$ has degree at most four, by the normal Klein four subgroup.
The elementary maximal-subgroup list $A_4,D_{10},S_3$ for $A_5$ gives
primitive degrees $5,6,10$. Comparing the specified $q+1$ values in
source Proposition 3.7 leaves only $q=9$ among its exceptional cases.

\section{Semilinear spread rigidity and the one-dimensional branch}
The following full argument is reused from the reviewed phase4 theorem
\emph{Rigidity of locally primitive semilinear spreads}. Its application
to the source question is already recorded there; this edition claims
no new question answer or priority. We retain Delandtsheer's reduction
quoted as source Proposition 5.2: a one-dimensional flag-transitive action
contains all translations, and its spread case consists of cosets of
equal-dimensional additive subspaces. The proof below includes the
restriction kernels and characteristic-dividing Galois case.

\begin{lemma}[Prime degree and arithmetic]\label{prime}
A finite primitive permutation group with cyclic normal subgroup and
cyclic quotient has prime degree, if its degree exceeds one. If
\[
 r=\frac{p^{u\ell}-1}{p^\ell-1}
\]
is prime, where $p$ is prime and $u>1$, then $u$ is prime,
$\ell=u^e$ for some $e\ge0$, and $p\not\equiv1\pmod u$.
\end{lemma}
\begin{proof}
Pass to the faithful induced action. The cyclic normal subgroup either
is trivial or is transitive. In the former case the whole induced group
is cyclic, hence regular of prime order. In the latter case the cyclic
normal subgroup is regular. Any proper nontrivial characteristic subgroup
would give a nontrivial block system preserved by the whole group, so
its order and the degree are prime.

If $u$ is composite, its geometric sum factors into two integer
geometric sums greater than one. Thus $u$ is prime. If a prime
$v\ne u$ divides $\ell$, put $x=p^{\ell/v}$. Reduction modulo
$\Phi_u(X)=1+X+\cdots+X^{u-1}$ shows that $\Phi_u(X)$ divides
$\Phi_u(X^v)$: multiplication by $v$ permutes the nonzero residues
modulo $u$. It follows that
$1<\Phi_u(x)<\Phi_u(x^v)=r$, contradicting primality. Hence
$\ell=u^e$. Finally $p\equiv1\pmod u$ would give $u\mid r$ while
$r>u$, again impossible.
\end{proof}

\begin{lemma}[Cyclic semilinear descent]\label{descent}
Let $p,u$ be primes with $p\not\equiv1\pmod u$. Let a finite group $J$
act faithfully and irreducibly on an $\F_p$-space $U$ of dimension $u^e$.
Suppose $C\triangleleft J$ is cyclic, $K=\F_p[C]$ is a field of
endomorphisms, $J$ normalizes $K$, and $J/C$ is a cyclic $u$-group.
Then $\dim_KU=1$.
\end{lemma}
\begin{proof}
Write $|K|=p^h$. Since $h\mid u^e$, $h$ is a power of $u$.
For every power $u^a$, the number $p^{u^a}-1$ is prime to $u$:
this is immediate for $p=u$, and otherwise follows from
$p^{u^a}\equiv p\not\equiv1\pmod u$. In particular $C$ has order
prime to $u$.

Choose a lift $t$ of a generator of $J/C$ and let $y$ be its
$u$-primary component inside the cyclic group generated by $t$.
The prime-to-$u$ component lies in $C$, so $J=C\langle y\rangle$.
Let $\sigma$ be the automorphism of $K$ induced by $y$ and let $m$
be its order. Then $z=y^m$ centralizes $K$, $C$ and $y$, hence $J$.
Every nonzero element of the commutative algebra $\F_p[z]$ is invertible:
its kernel is $J$-invariant, so irreducibility makes it injective. This
finite algebra is a field of order $p^a$, with $a\mid u^e$. Its
multiplicative group has order prime to $u$. Since $z$ has $u$-power
order, $z=1$. Thus the order of $y$ is exactly $m$.

There is a nonzero $y$-fixed vector without dividing by $m$. Choose
$0\ne v\in U$ and consider
\[
 S(b)=\sum_{i=0}^{m-1}\sigma^i(b)y^i(v),\qquad b\in K.
\]
Every $S(b)$ is fixed by $y$, by cyclically shifting the summands.
They cannot all vanish. Apply a $K$-linear functional nonzero on $v$;
vanishing would give a nonzero linear relation over $K$ among the
distinct field automorphisms $1,\sigma,\ldots,\sigma^{m-1}$.
Their linear independence follows by taking a shortest relation and
subtracting its value at $bx$ from a suitable multiple of its value
at $x$, shortening the relation. For some $b$ we therefore obtain a
nonzero fixed vector $w=S(b)$. Its $K$-line is $J$-invariant: $C$
acts by scalars and $y$ normalizes $K$ while fixing $w$.
Irreducibility gives $U=Kw$.
\end{proof}

\begin{theorem}[Locally primitive one-dimensional spread spaces]\label{spread}
Let $F=\F_{p^d}$, let $\mathcal S$ be a spread of $\ell$-dimensional
additive subspaces, with $0<\ell<d$, and let
$H\le\Gamma L_1(F)$ preserve it and act transitively. Suppose $H$
is primitive on $\mathcal S$ and $H_U$ is irreducible on a spread
member $U$. Then $\mathcal S$ is exactly the set of one-dimensional
$\F_{p^\ell}$-subspaces of $F$. Its size is prime and
$d=u^{e+1},\ell=u^e$ for a prime $u$.
Conversely such a locally primitive one-dimensional action exists
precisely when $(p^d-1)/(p^\ell-1)$ is prime; the geometry is then
uniquely $\operatorname{AG}(d/\ell,p^\ell)$ up to isomorphism.
This does not classify every possible acting subgroup.
\end{theorem}
\begin{proof}
The normal scalar subgroup of $H$ is cyclic and its quotient embeds
in the cyclic Galois group. Lemma~\ref{prime} gives prime spread size
$r=(p^d-1)/(p^\ell-1)$. Divisibility of the denominator into the
numerator implies $\ell\mid d$: dividing $d=a\ell+b$, $0\le b<\ell$,
would otherwise give $p^\ell-1\mid p^b-1$. Apply the arithmetic part
of that lemma with $u=d/\ell$.

Put $C_0=H_U\cap F^*$, and let $J$ be the faithful image of $H_U$
on $U$. The ambient scalar group $C_0$ restricts faithfully on nonzero
$U$, and $K=\F_p[C_0]$ is an actual subfield of $F$. Also $U$ is a
$K$-space. The semilinear stabilizer normalizes $K$; its faithful
image modulo the image $C$ of $C_0$ is a quotient of a subgroup of
$\operatorname{Gal}(F/\F_p)$, hence a cyclic $u$-group because
$d=u^{e+1}$. Lemma~\ref{descent} now gives $U=wK$ with
$|K|=p^\ell$. The unique subfield of that order is $\F_{p^\ell}$.
All $H$-images of $U$ are $K$-lines, since scalar-Frobenius maps
preserve this subfield. There are exactly $r$ such lines, so these
are the entire spread.

For the converse use $F^+$\,$\rtimes F^*$ on the $K$-line spread.
It is flag-transitive, its pencil action is the cyclic regular
prime-degree action of $F^*/K^*$, and its line stabilizer contains
$K^+\rtimes K^*$, sharply two-transitive on $K$. Its global action
on $F$ is sharply two-transitive as well. The local-action criterion
proves every required property. This establishes existence within a
one-dimensional group, without borrowing the full automorphism group
of the affine space.
\end{proof}

\begin{proposition}[Local two-transitivity in the one-dimensional branch]\label{onedim2t}
Under the spread hypotheses, local two-transitivity occurs only for
$\operatorname{AG}(2,2)$ with $A\Gamma L_1(4)$, or
$\operatorname{AG}(2,4)$ with $A\Gamma L_1(16)$.
\end{proposition}
\begin{proof}
Write $v=p^{u^{e+1}}$, $k=p^{u^e}$ and $d=u^{e+1}$ by
Theorem~\ref{spread}. Pencil two-transitivity forces
$r(r-1)\mid d(v-1)$. Since $v-1=r(k-1)$, we have
$r-1\mid d(k-1)$. But $k\mid r-1$ and $\gcd(k,k-1)=1$,
so $k\mid d$. Hence $p=u$ and $u^e\le e+1$, leaving only $e=0$
or $e=1$; in the latter case $p=u=2$. If $e=0$, cancellation gives
$1+p+\cdots+p^{p-2}\mid p-1$, which is possible only for $p=2$.
Thus the only parameter pairs are $(v,k)=(4,2),(16,4)$.

For four points, pencil two-transitivity requires a group order
divisible by $4\cdot3\cdot2=24$, so the full $A\Gamma L_1(4)$ is
necessary and is $S_4$. For sixteen points, the five-point pencil
requires $16\cdot5\cdot4\mid |G|$, and the twenty lines with four
points each require $20\cdot4\cdot3\mid |G|$. Their least common
multiple is $960$, exactly $|A\Gamma L_1(16)|$; the full group is
necessary. It suffices: on the five cosets $\F_{16}^*/\F_4^*$ the
multipliers are a regular cyclic group of order five and Frobenius
acts by multiplication by two, of order four on that quotient.
This is sharply two-transitive. On a line $\F_4$, its translations
and nonzero scalar multipliers give a sharply two-transitive group.
\end{proof}

\section{Scalar quotients and the circle branches}
\begin{lemma}[A scalar-orbit test]\label{scalarorbits}
Let $F$ be a finite field and $C\le F^*$. Suppose
$F^+\rtimes C\le H\le A\Gamma L_1(F)$. The affine action of $H$
is two-transitive exactly when its zero stabilizer is transitive on
$F^*/C$. If $\F_p[C]=F$, the action is primitive.
\end{lemma}
\begin{proof}
Translations reduce ordered-pair transitivity to transitivity of the zero
stabilizer on $F^*$. Scalar multiplication by $C$ is transitive on each
of its multiplicative cosets; these cosets are preserved as a system by
all scalar-Frobenius maps. This proves the first equivalence. A block
through zero is an additive subspace invariant under the zero stabilizer.
If it contains a nonzero $v$ and is invariant under $C$, it contains
$\F_p[C]v=F$. Thus no proper nonzero block exists.
\end{proof}

\begin{proposition}[Projective-line circles]\label{projectiveline}
Let $\operatorname{PSL}_2(q)\le G\le\operatorname{P\Gamma L}_2(q)$
in its natural degree-$q+1$ action. Its circle is always residually
weakly primitive. It is locally two-transitive exactly when the induced
action of $G$ on the determinant square classes is transitive. For even
$q$ this condition is automatic; for odd $q$ it says that some element
has nonsquare determinant in its semilinear projective representative.
\end{proposition}
\begin{proof}
The natural special-linear action is two-transitive. Fixing infinity
identifies its remaining points with $\F_q$. Its induced stabilizer
is the affine group $\F_q^+\rtimes C$, where
$C=(\F_q^*)^{\gcd(2,q-1)}$: the diagonal matrix
$\operatorname{diag}(a,a^{-1})$ gives multiplication by $a^2$,
and upper unipotents give every translation. General semilinear
stabilizers give every allowed scalar-Frobenius map. The quotient
character is determinant modulo squares, invariant under projective
scalar changes. Since the special-linear group is transitive, the full
point stabilizer maps onto the outer quotient. Lemma~\ref{scalarorbits}
therefore gives the claimed two-transitivity criterion on the remaining
points, equivalent to three-transitivity on the whole set.

For primitivity, if $q$ is prime, the additive group already has prime
degree. Otherwise every proper subfield has size at most $\sqrt q$,
whereas $|C|\ge(q-1)/2>\sqrt q-1$. Thus $C$ cannot be contained in a
proper subfield and generates $\F_q$; Lemma~\ref{scalarorbits} applies.
Together with Proposition~\ref{circles} this proves the assertion.
The criterion retains mixed diagonal-field extensions; it does not
incorrectly require the acting group to contain $\operatorname{PGL}_2(q)$.
\end{proof}

\begin{lemma}[A characteristic subgroup in a regular normal subgroup]
If a transitive group on $X$ contains a regular normal subgroup $U$
with a nontrivial proper characteristic subgroup $Z$, then it is
imprimitive.
\end{lemma}
\begin{proof}
The $Z$-orbits are blocks, since $Z$ is normal in the full group.
Regularity identifies their size with $|Z|$, strictly between one and
$|X|=|U|$.
\end{proof}

This proves the local exclusion for the natural Suzuki and Ree circle
actions from their precise retained rank-one structure: a point
stabilizer has a regular unipotent radical of order $q^2$ or $q^3$,
respectively, and a proper nontrivial characteristic center. Field and
torus automorphisms preserve this center. Thus the remaining-point
action of the point stabilizer is imprimitive, also for its allowed
outer extensions. No assertion that a Suzuki ovoid is itself a thick
linear space is needed. Natural unitary circles instead preserve the
actual thick unital, so Proposition~\ref{circles} directly excludes them.
Projective actions of dimension at least two preserve their thick
projective lines and are excluded in the same way.

For affine circles over an odd prime field, the actual affine lines
are thick and again give an invariant pencil partition. For extension
field actions on an even field larger than two, the same exclusion
uses the actual invariant extension-field lines, after establishing
their normalizer as in Lemma~\ref{field}. Over $\F_2$, by contrast,
affine lines have only two points and cannot be used for this exclusion.
The natural special-linear family gives three-transitivity, and the
symplectic family gives primitivity but not two-transitivity on the
nonzero vectors by Proposition~\ref{symplectic}. The exceptional
$G_2(2)'$ degree-63 local action requires its own small case: it must
not be excluded by pretending $\operatorname{AG}(6,2)$ is thick.
The following Proposition~\ref{binaryg2} handles it using the natural
module as a retained input and a complete small maximal-subgroup
calculation; the other exceptional action interfaces are kept separate.

The one-dimensional circle branch is covered by Theorem~\ref{spread}
with $\ell=1$: the pencil size $p^d-1$ must be prime. This includes,
for example, $\operatorname{AGL}_1(8)$ with primitive cyclic pencil
of degree seven. It is omitted from the literal stand-alone list in
source Proposition 4.3, whose proof defers the one-dimensional case
to Section 5. The source's final Theorem 6.1 does include it in its
first branch. The corresponding locally two-transitive circle is
only the four-point case among one-dimensional spread actions.

For the exceptional almost-simple branches we will retain exact
permutation-action information, rather than infer primitivity from
abstract group names or orders alone. In particular the $A_7$ action
on fifteen points preserves the lines of $\operatorname{PG}(3,2)$.
It is two-transitive on points but its point stabilizer preserves
seven pairs on the other fourteen points. Its circle is not locally
primitive and cannot be three-transitive. This is the previously
accepted Topic G circle correction, not a new conjecture result.

\begin{proposition}[The binary exceptional local action]\label{binaryg2}
Retain the natural transitive binary representation of dimension six of
$N=G_2(2)'\cong\operatorname{PSU}_3(3)$. Every group
$N\le H\le N_{\GL_6(2)}(N)$ is primitive but not two-transitive on
the sixty-three nonzero vectors. Thus $2^6\rtimes H$ has a residually
weakly primitive circle action and is not locally two-transitive.
\end{proposition}
\begin{proof}
The small exact finite premise supplied by \texttt{small\_g2\_premise.g}
is the complete maximal-subgroup order list
\[
 |N|=6048,\qquad 96,96,168,216,\qquad |\Aut(N)|=12096.
\]
It is computed in GAP from $\operatorname{PSU}(3,3)$, including a
simplicity check and both actual degree-63 coset actions. This is a
required finite premise, not an assertion derived solely from a subgroup
order. In the natural transitive action a vector stabilizer has order
$6048/63=96$. Any proper overgroup of it lies in a maximal subgroup,
and no order in the displayed complete list is a larger multiple of
96. Thus the stabilizer is maximal and the action is primitive.

For the normalizer bound we can avoid assuming an unverified full
matrix normalizer. Transitivity on nonzero vectors implies irreducibility.
Its commuting algebra is a finite division ring, hence a field $\F_{2^e}$,
with $e\mid6$. If $e=3$ or $6$, Lagrange's theorem would embed $N$
in $\GL_2(8)$ or $\GL_1(64)$, whose orders are not divisible by
6048. If $e=2$, simplicity and perfectness put $N$ inside
$\operatorname{SL}_3(4)$; its intersection with the central subgroup of order three
is trivial, so it embeds in $\operatorname{PSL}_3(4)$, of order 20160, again not
divisible by 6048. Consequently $e=1$, and its invertible commuting
algebra is trivial. Conjugation embeds the entire linear normalizer
in $\Aut(N)$. A two-transitive action on sixty-three points would
require order divisible by $63\cdot62$, which does not divide 12096.
Every intermediate group inherits primitivity from $N$, but none
can be two-transitive. Proposition~\ref{circles} proves the conclusion.
\end{proof}

\begin{proposition}[Orthogonal-point symplectic actions]\label{orthogonalcircles}
Retain the standard two-transitive actions of
$\operatorname{Sp}_{2n}(2)$ of degrees
$2^{2n-1}\pm2^{n-1}$, $n\geq3$, and the classical rank-three subdegrees
of their orthogonal point stabilizers. Both circle actions are locally
primitive and neither is locally two-transitive.
\end{proposition}
\begin{proof}
Put $b=2^{2n-2}$ and $d=2^{2n-1}\pm2^{n-1}-1$.
On the $d$ remaining points, a point stabilizer of the orthogonal group
has orbits of lengths $1$, $a=d-1-b$, and $b$; these are the precise
retained rank-three data. A block containing the fixed point is a union
of these orbits, so its only possible proper nontrivial sizes are
$a+1=d-b$ and $b+1$. The first is an odd integer greater than one; if it
divided $d$, it would divide the power of two $b$, impossible. For the
minus sign, $d/2<b+1<d$. For the plus sign,
$2(b+1)<d<3(b+1)$, since $n\geq3$. Thus $b+1$ cannot divide $d$ either.
The local action is primitive. Its rank is three, so it is not
two-transitive. The circle criterion proves both conclusions.
\end{proof}

\section{The sharply flag-transitive and nonspread branches}
We retain the projective flag-transitivity classification in source
Proposition 3.1, attributed to Kantor and Cameron--Kantor: besides the
classical group and the $A_7$ action on $\operatorname{PG}(3,2)$, a
flag-transitive projective plane has prime point number $m^2+m+1$ and
a sharply flag-transitive Frobenius group of order
$(m^2+m+1)(m+1)$. We also retain Feit's restrictions in Proposition 3.2:
the Desarguesian possibilities are $m=2,8$, while another example
requires $m\equiv0\pmod8$ and
$d^{m+1}\equiv1\pmod{m^2+m+1}$ for every divisor $d$ of $m$.
The reported computer extension of Feit's bound to $m\le10^{10}$ is
not a premise of the following proof and has no supplied source program.

\begin{proposition}[The exceptional projective branch]\label{projectivefrobenius}
A locally primitive sharply flag-transitive projective plane in this
retained branch is the Fano plane with group $7:3$. It is not locally
two-transitive.
\end{proposition}
\begin{proof}
The point stabilizer acts regularly on its $m+1$ incident lines.
A regular action is primitive only in prime degree, so $m+1$ is
prime. For a putative non-Desarguesian example, put $N=m^2+m+1$.
We have $m^3\equiv1\pmod N$, and the order of $m$ modulo the prime
$N$ is exactly three, since $1<m<N$. Feit's condition with $d=m$
therefore gives $3\mid m+1$, forcing $m=2$, a contradiction.
Among the Desarguesian possibilities, only $m=2$ has $m+1$ prime.
The group $7:3$ is regular on flags of the Fano plane; both local
three-point groups are cyclic regular. They are primitive but are
not two-transitive.
\end{proof}

The other exceptional projective action has $A_7$ on fifteen points
and thirty-five lines of $\operatorname{PG}(3,2)$. Retain its classical
identification and point stabilizer $\operatorname{PSL}_3(2)$ from
Cameron--Kantor/Taylor, as in source Proposition 3.3. The pencil is its
natural seven-point plane action, two-transitive by the elementary
basis argument. The stabilizer of a line induces $S_3$ on its three
points. With these exact retained local identifications, its claimed
local two-transitivity follows. This geometric action must not be
confused with the circle on the same fifteen points.

\begin{proposition}[One-dimensional nonspread actions, including the small boundary]\label{nonspread}
Retain Delandtsheer's translation/nonspread reduction quoted in source
Propositions 5.2 and 5.7. A flag-transitive one-dimensional nonspread
linear space is locally primitive precisely in the following two cases:
the three-point circle with $S_3=\operatorname{AGL}_1(3)$, and the
Fano plane with $7:3$. Only the first is locally two-transitive.
\end{proposition}
\begin{proof}
Write $v=p^d$ and let $k$ and $r=(v-1)/(k-1)$ be line size and pencil
size. The point stabilizer is a subgroup of $\Gamma L_1(p^d)$,
cyclic by cyclic. Its transitive action is primitive exactly when
$r$ is prime, by Lemma~\ref{prime} and the elementary fact that a
transitive prime-degree action is primitive. The retained nonspread
reduction gives $v\mid b$, where $b=vr/k$, and hence $k\mid r$.
Thus $k=r$ and $v=k(k-1)+1$.

There are now $b=v$ lines. Every other line meets a fixed line $L$:
its $k$ points lie on $k-1$ other lines each, giving $k(k-1)=v-1$
distinct lines, since two distinct points of $L$ cannot share another
line. If $k=2$ this is exactly the three-point circle. Flag transitivity
requires all six permutations, so the specified group is $S_3$; both
local actions have degree two.

If $k\ge3$, this is a projective plane of order $k-1$. The retained
projective classification and the preceding proposition leave exactly
the Fano plane and its group $7:3$. Conversely these two structures
have the asserted one-dimensional realizations and primitive local
actions. The cyclic local groups in the Fano case are not two-transitive.
\end{proof}

The source's Proposition 5.8 omits the $k=2$ case when passing from
$r=k$ to a nondegenerate projective plane. Its earlier Proposition 5.1
already lists this three-point circle. The footnote after Proposition 5.6
saying the spread restriction can be removed must therefore include this
extra three-point action. The final circle branches of Theorems 6.1
and 6.2 already include it as the natural degree-three $\operatorname{PSL}_2(2)$
action. No new family beyond that final union is being claimed.

\section{Comparing different affine structures}
The printed proof of source Lemma~3.5 assumes that two invariant affine
structures are nested merely because their line sizes are ordered. The
following comparison does not make that assumption.

\begin{lemma}[An intrinsic translation group]\label{intrinsictranslations}
Let $G$ act primitively on a set $X$ of size $p^d$, $d\geq2$, and contain
a regular elementary abelian normal subgroup $V$. If $G$ preserves a
Desarguesian affine-space structure on $X$, then $V$ is its translation
group.
\end{lemma}
\begin{proof}
The point stabilizer acts irreducibly on $V$: a proper invariant subgroup
would give a nontrivial block system of its cosets. Thus $V$ is minimal
normal in $G$. In the given affine structure, project $G$ to its cyclic
field-automorphism group. The kernel of the restriction to $V$ is normal
in $G$ and is nontrivial, because an elementary abelian group of rank at
least two cannot inject into a cyclic group. Minimality puts all of $V$
in the affine linear subgroup.

Take linear parts there. A finite $p$-subgroup of a linear group in
characteristic $p$ has a nonzero fixed vector: orbit counting on the
vector space makes the number of fixed vectors divisible by $p$.
The common fixed space $W$ of the image of $V$ is invariant under the
entire semilinear image of $G$, since $V\lhd G$. If $W$ were proper,
its affine cosets would be a nontrivial $G$-invariant partition of $X$.
Thus $W$ is the whole vector space. The linear image of $V$ is trivial,
so $V$ consists of translations; equality follows from the common order
$p^d$.
\end{proof}

\begin{proposition}[Scalar rigidity]\label{scalarrigidity}
In Lemma~\ref{intrinsictranslations}, identify $X$ with $V$ by choosing
an origin. Suppose its stabilizer contains a perfect subgroup $N$ and
$\End_{\F_p N}(V)=\F_q$. Then every $G$-invariant Desarguesian affine
structure has line size at most $q$. Its scalar field is an actual
subfield of this commuting algebra. If the natural $\F_q$-structure is
$G$-invariant, it has largest line size among these structures.
\end{proposition}
\begin{proof}
All structures have the same translation group by the lemma, hence the
same addition on $V$. In another structure over $\F_Q$, the field action
of $N$ is a homomorphism to a cyclic group and is trivial because $N$ is
perfect. Thus $N$ is $\F_Q$-linear. Every scalar of $\F_Q$ commutes with
$N$, giving an embedding of that actual scalar field into
$\End_{\F_p N}(V)=\F_q$. In particular $Q\leq q$. This argument
constructs the common scalar setting rather than assuming a prior
partition of one pencil by the other.
\end{proof}

For natural special-linear modules of dimension at least three and natural
symplectic modules of dimension at least four in the thick range $q\geq3$,
the perfectness and the commuting-algebra calculations above apply.
For the even-field natural $G_2(q)$ module, the commuting algebra follows
from an order argument. Retain its standard faithful six-dimensional
representation, transitivity on nonzero vectors, simplicity for $q\geq4$,
and order $q^6(q^6-1)(q^2-1)$. Transitivity gives irreducibility over
$\F_p$. Its commuting algebra is a finite division ring, hence a field
containing $\F_q$, of relative degree dividing six. Relative degrees three
or six would embed the group into $\GL_2(q^3)$ or $\GL_1(q^6)$, impossible
by the $p$-parts of their orders. Relative degree two would embed this
perfect simple group into $\operatorname{SL}_3(q^2)$ and then, with
trivial central intersection, into $\operatorname{PSL}_3(q^2)$.
The ratio of the latter order to $q^6(q^6-1)(q^2-1)$ is
$(q^2+1)/3$, not an integer since $q$ is a power of two. Thus the
commuting algebra is exactly $\F_q$. The $q=2$ derived-group case
has been proved separately in Proposition~\ref{binaryg2}.

For the one-dimensional branch no comparison of affine structures is
needed to bound line size. A primitive pencil has prime size by
Lemma~\ref{prime}; if its line size is $p^f$ and point size $p^d$, that
lemma gives $d=u^{a+1}$ and $f=u^a$ for a prime $u$. Every nondegenerate
affine space on $p^d$ points has line size $p^b$ where $b$ is a proper
divisor of $d$, so $b\leq u^a=f$. This proves maximality even against
non-Desarguesian affine planes.

\begin{proposition}[The natural even-field exceptional pencil]\label{naturalg2}
Retain the classical six-dimensional symplectic representation of
$G_2(q)$, $q\geq4$ even, in which the points of the split Cayley hexagon
are exactly the vector lines, and their stabilizers are maximal
parabolic subgroups. Every group between this group and its full
$\F_2$-linear normalizer is primitive but not two-transitive on vector
lines. With translations its affine action is locally primitive but not
locally two-transitive. The analogous binary derived-group assertion
is Proposition~\ref{binaryg2}.
\end{proposition}
\begin{proof}
Transitivity and maximality of the stated point stabilizers imply
primitivity; these are the exact retained building inputs, attributed
to the classical hexagon identification used in source Proposition~3.6.
The commuting-field calculation above makes every normalizer element
semilinear over $\F_q$. Remove its field automorphism and pull back the
preserved nondegenerate alternating form. The resulting form is again
$G_2(q)$-invariant. Its defining operator relative to the original form
commutes with $G_2(q)$, so it is a scalar by the calculated commuting
algebra. Hence the entire normalizer preserves perpendicularity on
vector lines. There are distinct perpendicular and nonperpendicular
pairs in dimension six, so none of these groups is two-transitive on
the pencil. The retained transitivity of the natural group on nonzero
vectors implies transitivity of a line stabilizer on its nonzero
vectors. Translations then give two-transitivity on each affine line.
Proposition~\ref{translation} proves the remaining assertions.
\end{proof}

\begin{lemma}[Large vector-stabilizer orbits]\label{largeorbits}
Let $V=\F_q^n$, $q\geq3$. If $n\geq3$ and
$N=\operatorname{SL}_n(q)$, then $N_v$, for $0\ne v\in V$, is
transitive on $V\setminus\F_qv$, of size $q^n-q>q^{n/2}$.
If $n\geq4$ is even and $N=\operatorname{Sp}_n(q)$, every
$N_v$-orbit outside $\F_qv$ has size at least
$q^{n-1}-q>q^{n/2}$.
Consequently a $G$-invariant affine-space structure on $V$, where
$G_0\geq N$, has no line larger than $q$.
\end{lemma}
\begin{proof}
For the special-linear assertion, extend $(v,w)$ and $(v,w')$ to bases
for any $w,w'\notin\F_qv$. Map the bases to each other and correct
the determinant on a third basis vector. This fixes $v$ and sends $w$
to $w'$. For the symplectic assertion, symplectic basis extension gives
one orbit $v^\perp\setminus\F_qv$, of size $q^{n-1}-q$, and one
orbit $\{w:\langle v,w\rangle=c\}$, of size $q^{n-1}$, for each
$c\ne0$. The first transitivity follows by extending an ordered
independent isotropic pair; the second follows by extending the
hyperbolic pair $(v,c^{-1}w)$. These constructions account for every
vector outside $\F_qv$. The displayed strict inequalities hold already
at $n=3$ and $n=4$, respectively, and increase thereafter.

Here the comparison is only with affine spaces, not arbitrary linear
spaces. Every nondegenerate finite affine space on $|V|=q^n$ points
has line size at most $\sqrt{|V|}$: a plane has exactly $k^2$ points
if its line size is $k$, and higher dimension gives at least a plane.
For an alternative line $L$ through $0$ and $v$, uniqueness of the line
through these two points makes $L$ invariant under $N_v$.
It cannot contain any of the above outside orbits. Thus
$L\subseteq\F_qv$, and its size is at most $q$.
\end{proof}

\begin{theorem}[Maximal line size for a primitive affine pencil]
\label{affinemax}
Let a fixed faithful permutation group $G$ on a finite set $X$ act
flag-transitively on an affine space $\Gamma$ with at least two lines
through each point. If $G_x$ is primitive on its pencil, then
$\Gamma$ has largest line size among all $G$-invariant affine-space
structures on $X$ on which $G$ is flag-transitive.
This recovers source Lemma~3.5 by contraposition. Its proof retains the
precise BDDKLS--Huybrechts affine classification and the BDD exclusion
of local point-primitivity for the full L\"uneburg-plane group.
\end{theorem}
\begin{proof}
A plane already has line size $\sqrt{|X|}$, the universal affine bound
just proved. If the lines of $\Gamma$ have size two, a larger-line
structure would partition each pencil of the circle into the sets
$L'\setminus\{x\}$, contradicting Proposition~\ref{circles}.
We may therefore suppose $\Gamma=\operatorname{AG}(n,q)$,
$n\geq3$, $q\geq3$; the coordinatization of affine spaces of dimension
at least three is the classical retained input.

Use exactly the classification interface quoted in source Section~3.2.1:
a flag-transitive group is one-dimensional, belongs to the excluded
$\operatorname{AG}(4,3)$ last-derived-$2.A_5$ case, or is two-transitive
with one of Huybrechts' six stated natural representation types.
The excluded case is a retained impossibility, not inferred here from
its abstract group name. In the one-dimensional case the prime-pencil
and prime-power argument immediately preceding this theorem proves the
line-size bound, including comparison with non-Desarguesian planes.

In the two-transitive cases write $G=V\rtimes G_0$. The three small
exceptional representation types have imprimitive pencils by the
explicit exclusions in the classical-pencils section. For a natural
special-linear, symplectic or even-field $G_2$ normal subgroup over
$\F_{q^t}$, the commuting-algebra arguments make $G_0$ normalize that
actual scalar field. If $t>1$, its lines partition the $\F_q$-pencil
nontrivially, contrary to Lemma~\ref{field}. (The special-linear
one-dimensional case belongs to the one-dimensional alternative.)
Thus the remaining modules are natural over $\F_q$ itself. In the
special-linear and symplectic cases Lemma~\ref{largeorbits} applies
also to non-Desarguesian alternative planes and proves the conclusion.

It remains to consider $N=G_2(q)$ on $V=\F_q^6$, with $q$ even and
$q\geq4$. This subgroup is perfect and its commuting field is
$\F_q$, as proved above. Scalar rigidity bounds every alternative
Desarguesian line size by $q$. Any remaining alternative is a
non-Desarguesian affine plane. It cannot have a one-dimensional
flag-transitive group: every subgroup of $A\Gamma L_1(p^d)$ is
soluble, whereas $G$ contains the nonabelian simple subgroup $N$.
The retained BDD list of non-one-dimensional flag-transitive
non-Desarguesian planes consists of the L\"uneburg planes, the nearfield
plane of order nine and the Hering plane of order twenty-seven.
The latter two have an odd number of points, whereas $|V|=q^6$ is
even. Thus only a L\"uneburg plane could remain.

But $G$ is two-transitive on $V$, independently of its incidence
structure. In any invariant linear space this implies two-transitivity
of every line stabilizer on that line: an element sending one ordered
pair of distinct points of the line to another sends their unique
joining line to itself. Hence that line action is primitive, and remains
primitive in the full automorphism group by Lemma~\ref{coarsening}.
This contradicts the retained BDD fact, quoted in source Section~3.2.2,
that the full L\"uneburg-plane group is not locally point-primitive.
This eliminates the last alternative and completes the proof.
\end{proof}

The argument does not assert that two arbitrary invariant affine
structures are nested. Its orbit bounds and classification comparison
replace the unproved nesting sentence and the unnecessary full-group
equality in the printed proof. The classification inputs above remain
charged retained inputs; their proofs have not been reconstructed here.


\section{The exact scalar obstruction in a Hermitian pencil}
Let $q=p^f\geq2$, $F=\F_{q^2}$, and put $d=\gcd(3,q+1)$.
The classical unital
has the isotropic projective points for
\[
 h(X)=X_0^qX_2+X_0X_2^q+X_1^{q+1}.
\]
Its lines are the secants, with incidence inherited from the plane.
The standard full projective semilinear unitary group and its special
subgroup are denoted by $\PGammaU_3(q)$ and $\PSU_3(q)$.
We retain the classical identification of the full unital automorphism
group with $\PGammaU_3(q)$. The local computations below are explicit and
do not retain the blanket local-two-transitivity assertion of the source.

Let $C=(F^*)^d$ and $\mathcal C=F^*/C$. There is a natural action
\[
 \pi:\PGammaU_3(q)/\PSU_3(q)\longrightarrow\operatorname{Sym}(\mathcal C).
 \tag{U1}
\]
To specify it, fix $\infty=\langle(1,0,0)\rangle$ and identify its pencil
with $F$. A point stabilizer acts as $b\mapsto ab^{p^i}+c$; send this to
$Cz\mapsto Ca z^{p^i}$. Changing a representative by a special-unitary
stabilizer multiplies $a$ by $C$ and changes only the translation $c$.
Because the special-unitary group is point-transitive, the full group's point stabilizer
surjects onto the full outer quotient. This defines (U1). When $d=3$,
choosing a generator of $F^*/C$ identifies the action with
$j\mapsto t+p^i j$ on $\mathbb Z/3\mathbb Z$.

\begin{theorem}[Unitary local criterion]\label{unitary}
Suppose $\PSU_3(q)\leq G\leq\PGammaU_3(q)$ in its specified unital action.
Then $G$ is flag-transitive and its line stabilizers are two-transitive on
their incident points. The point stabilizers are two-transitive on their
incident lines if and only if $\pi(G)$ is transitive on $\mathcal C$.
For $q>2$ the action is always residually weakly primitive. For $q=2$
it is residually weakly primitive if and only if the same transitivity
condition holds; in that case it is locally two-transitive.
\end{theorem}
\begin{proof}
Every unital point other than $\infty$ has a unique representative
$(a,b,1)$ with $a+a^q+b^{q+1}=0$. Its joining line with $\infty$ depends
only on $b$. Consequently the pencil at $\infty$ identifies with $F$.
For $c,t\in F$ satisfying $t+t^q=-c^{q+1}$, the matrices
\[
 u(c,t)=\begin{pmatrix}1&-c^q&t\\0&1&c\\0&0&1\end{pmatrix}
 \tag{U2}
\]
preserve $h$ and have determinant one. Their action on the pencil is
$b\mapsto b+c$. The induced translation group therefore is all $F^+$.
The determinant-one diagonal matrices
\[
 \operatorname{diag}(z,z^{q-1},z^{-q}),\quad z\in F^*,
 \tag{U3}
\]
induce $b\mapsto z^{2q-1}b$. The image of this power map is $C$, because
$\gcd(2q-1,q^2-1)=\gcd(3,q+1)=d$.

Conversely, a projective unitary element fixing $\infty$ preserves its
orthogonal tangent. In the displayed basis its matrix is upper triangular.
The unitary equations give (U2) followed by a unitary diagonal matrix;
the determinant-one condition gives (U3), up to a projective scalar.
Thus the induced special-unitary pencil group is exactly
$F^+\rtimes C$, not merely a subgroup of that group. General projective
unitary diagonals give all multipliers in $F^*$, and field automorphisms
give $b\mapsto b^{p^i}$. This also verifies the outer action (U1).

The full group $G_\infty$ contains every translation. Removing a
translation from each element fixing an ordered pair reduces
pencil two-transitivity to transitivity on $F^*$ of its zero stabilizer.
That stabilizer contains $C$, which acts regularly on each multiplicative
coset of $C$. It is therefore transitive on $F^*$ exactly when the induced
action on $\mathcal C$ is transitive. This proves the claimed criterion.

For primitivity, $C$ generates the field $F$ over $\F_p$ when $q>2$.
Indeed $|C|=(q^2-1)/d>q-1$, whereas the multiplicative group of any proper
subfield of $F$ has order at most $q-1$. Thus no nonzero proper additive
subspace can be invariant under $C$. The affine pencil action is
primitive by Proposition~\ref{translation}'s block argument. When $q=2$,
$C=1$, and the zero stabilizer lies in $\Gamma L_1(4)=\GL_2(2)$.
Its action on the four-point additive space is irreducible exactly when
it contains an element of order three, equivalently when it is transitive
on the three nonzero vectors. This is precisely the criterion on
$\mathcal C$ in this case.

For completeness the line action is two-transitive already in the special
unitary subgroup. A secant spans a nondegenerate two-dimensional Hermitian
space and its orthogonal complement is a nonsingular point. The subgroup
acting as $\operatorname{SU}_2(q)$ on the former and identically on the
latter lies in $\operatorname{SU}_3(q)$. In a hyperbolic basis the isotropic
points of the former are $\infty$ and $\langle(t,1)\rangle$ with
$t+t^q=0$. The matrices $\left(\begin{smallmatrix}1&c\\0&1\end{smallmatrix}\right)$,
$c+c^q=0$, fix $\infty$ and act regularly on the other $q$ points.
For any nonzero such $c$, the determinant-one isometry
$\left(\begin{smallmatrix}0&c\\-c^{-1}&0\end{smallmatrix}\right)$ interchanges
$\infty$ and zero. Hence this group is two-transitive on all $q+1$ points.

The same rank-one isometries give point-transitivity of $\PSU_3(q)$,
and (U2) gives transitivity on each pencil. Thus it is flag-transitive.
Global point-primitivity follows, for example, from its usual two-transitive
unital action: (U2) is regular on the $q^3$ other unital points. Together
with the pencil and line calculations this proves every assertion.
\end{proof}

\begin{corollary}[Two literal source boundaries]
For $d=1$, every group between $\PSU_3(q)$ and $\PGammaU_3(q)$ is locally
two-transitive. For $d=3$, $\PSU_3(q)$ itself is not locally two-transitive.
In particular $\PSU_3(2)$ is not residually weakly primitive on the unital
$\operatorname{AG}(2,3)$, although $\PGU_3(2)$ is.
\end{corollary}
\begin{proof}
If $d=1$ there is one coset. If $d=3$, the special-unitary image in (U1)
is trivial on three cosets. The last assertion uses the $q=2$ part of
Theorem~\ref{unitary}; the projective unitary image supplies the cyclic
regular action on those three cosets. As a small order check,
$|\PSU_3(2)|=72$ and its point stabilizer has order eight, so it cannot
be two-transitive on the four incident lines.
\end{proof}

\section{Witt--Bose--Shrikhande local actions}
We use the standard nonsplit-torus construction of the
Witt--Bose--Shrikhande space, attributed in source Section~3.4 to the earlier
work cited there. Its full automorphism group
$\operatorname{P\Gamma L}_2(q)$ is the retained classification input from
BDD, as used by the source; the local-action argument below is given in
full. In particular, the full-group identification is not inferred from
an exhibited subgroup. Since $q$ is even,
$\operatorname{PSL}_2(q)=\operatorname{PGL}_2(q)=\operatorname{SL}_2(q)$.

\begin{proposition}[Torus--involution incidence]\label{wbs-model}
Let $q=2^e$, $e\geq3$, and $G=\operatorname{SL}_2(q)$. Let the points be
normalizers $D=N_G(T)$ of nonsplit tori $T$ of order $q+1$, and let the lines
be the involutions $s$ of $G$. Declare $D$ incident with $s$ when $s\in D$.
This is a flag-transitive linear space with
\[
 v=q(q-1)/2,\qquad b=q^2-1,\qquad k=q/2,\qquad r=q+1.
\]
A point stabilizer in $G$ is $D$ of order $2(q+1)$, a line stabilizer is
$U=C_G(s)$ of order $q$, and a flag stabilizer is $\langle s\rangle$.
\end{proposition}
\begin{proof}
Identify the two-dimensional vector space with $\F_{q^2}$. Multiplication
by its norm-one elements gives the cyclic group $T$. The map $z\mapsto z^q$
inverts $T$ and, in characteristic two, has determinant one. The normalizer
is therefore the dihedral group $D$ of order $2(q+1)$: its elements either
fix or exchange the two common eigenlines over $\F_{q^2}$. Conversely these
are all possibilities for a normalizer. Every nonidentity element of $T$
has distinct eigenvalues outside $\F_q$, since
$\gcd(q-1,q+1)=1$, and its centralizer in $G$ is exactly $T$.
All nonsplit tori are conjugate: conjugate their two-dimensional field
embeddings, and adjust a conjugating matrix's determinant by a centralizing
field multiplier, whose norm ranges over $\F_q^\times$. Also $D$ is
self-normalizing, since its odd cyclic subgroup is characteristic.
This proves the point count.

Every involution is conjugate to
$s=\left(\begin{smallmatrix}1&1\\0&1\end{smallmatrix}\right)$.
Its centralizer consists exactly of
$u_a=\left(\begin{smallmatrix}1&a\\0&1\end{smallmatrix}\right)$,
$a\in\F_q$, as direct matrix multiplication shows. Thus $U$ is elementary
abelian of order $q$ and there are $q^2-1$ involutions. The $q+1$ reflections
of $D$ form one $T$-orbit, because squaring is bijective on $T$.
Consequently $G$ is flag-transitive and the flag stabilizer is
$C_D(s)=\langle s\rangle$. This proves $r=q+1$ and $k=q/2$.

Two distinct point normalizers cannot contain two distinct common
involutions: their nontrivial product lies in the rotation torus of each,
and its unique centralizer identifies those tori and hence the normalizers.
Thus two points lie on at most one line. Moreover
\[
 b\binom{k}{2}=\frac{(q^2-1)q(q-2)}8
             =\binom{q(q-1)/2}{2}.
\]
Every pair therefore lies on exactly one line. The parameter bounds give
at least two points on every line and two lines on every point.
\end{proof}

\begin{theorem}[Uniform local exclusion and a corrected source sentence]
\label{wbs-exclusion}
For $e\geq3$, no flag-transitive subgroup of the full group
$A=\operatorname{P\Gamma L}_2(2^e)$ of $W(2^e)$ is residually weakly
primitive. Nevertheless the full group of $W(8)$ is primitive on the points
of each line: its induced line action is $\operatorname{AGL}_1(4)\cong A_4$.
Thus the final WBS exclusion of source Section~3.4 is valid, but its stronger
assertion that the full group is never locally point-primitive is false.
\end{theorem}
\begin{proof}
First consider the pencil at $D$. Its lines are the $q+1$ involutions in
$D$. The rotation subgroup $T$ acts regularly on them. It is normal in
$A_D$, since it is characteristic in $D$ and $G\lhd A$. If $q+1$ is
composite, a proper nontrivial characteristic subgroup of the cyclic $T$
gives a nontrivial $A_D$-invariant partition into its orbits. Thus pencil
primitivity forces $2^e+1$ to be prime. Writing $e=2^a m$ with $m$ odd
shows that $m=1$: for $m>1$, $2^{2^a}+1$ is a proper divisor. Hence
$e$ is a power of two, and $e\geq4$ in the present range.

Now consider the line represented by the displayed involution $s=u_1$.
Its stabilizer is the centralizer $A_s$. Each element of $A$ is a matrix
part followed by a field automorphism. Field automorphisms fix $s$;
the matrix centralizer calculation above gives
\[
 A_s=U\rtimes\langle\sigma\rangle,
 \qquad \sigma(a)=a^2,\qquad |\langle\sigma\rangle|=e.
\]
Fix a point $D$ on this line and put $B=A_s\cap A_D$.
Since $G$ is flag-transitive, every coset of $G$ in $A$ has a representative
fixing this flag. Hence $B/(B\cap G)\cong A/G\cong C_e$, whereas
$B\cap U=B\cap G=\langle s\rangle$. Also $A_s=UB$.
The line points can consequently be identified with the additive quotient
\[
 V=U/\langle s\rangle\cong\F_{2^e}/\F_2,
\]
with $V$ acting regularly. Conjugation by the point stabilizer acts on this
quotient through $\langle\sigma\rangle$; translation parts have no effect
on that action. The affine block criterion in
Proposition~\ref{translation} says that line primitivity is equivalent to
irreducibility of this stabilizer on $V$.

For the only remaining pencil possibility, $e$ is a power of two. The
stabilizer image on $V$ is therefore a $2$-group. A $2$-group acting on a
nonzero vector space over $\F_2$ has a nonzero fixed vector: all its nonfixed
orbits have even size, so the number of fixed vectors is even and includes
zero. Since $\dim_{\F_2}V=e-1>1$, a fixed one-dimensional subspace is a
proper invariant subspace. The line action is imprimitive. Thus the full
group fails at least one of the two required local actions for every $e$.
A primitive local action of a subgroup would make the corresponding full
group action primitive by Lemma~\ref{coarsening}; no subgroup can evade
this obstruction. This proves the uniform exclusion.

For the stronger source assertion, set $e=3$. On $V=\F_8/\F_2$, a class
fixed by $\sigma$ has $x^2+x\in\F_2$. The value zero gives the zero class;
the value one would put a root of $X^2+X+1$ in $\F_8$, impossible because
its roots have multiplicative order three whereas $|\F_8^\times|=7$.
The order-three action on the two-dimensional $V$ is therefore irreducible.
It is transitive on the three nonzero vectors. The induced line group is
$V\rtimes C_3$ of order twelve, proving the claimed primitive action.
Its pencil remains imprimitive since $q+1=9$. The short accompanying
$q=8$ calculation validates these exact stabilizers but is not a premise
of the written proof.
\end{proof}


\section{Dependency and review boundary}
The companion application supplies the full source dispositions and both
corrected classification converses. All proofs in this common paper are
written here; no earlier common theory is treated as a free imported proof.
The semilinear spread descent is reused from the already reviewed phase-four
argument with its prior scope unchanged. The large classification inputs
and classical representation identifications are retained at their named
point of use. The binary exceptional pencil uses a complete small finite
premise, whereas the WBS diagnostic only validates a uniform written proof.
No new Lean declaration or whole-source admission is claimed by this private
unreviewed edition.

\begin{thebibliography}{1}
\bibitem{BDL} F. Buekenhout, P.-O. Dehaye and D. Leemans,
\emph{RWPri and $(2T)_1$ flag-transitive linear spaces},
Beitr\"age zur Algebra und Geometrie 44 (2003), 25--46.
\end{thebibliography}
\end{document}
